problem:  
Let \( (M^n, g) \) be a complete, noncompact Riemannian manifold with **nonnegative sectional curvature**.  
For a fixed base point \( p \in M \), let  
\[
r(x) = \mathrm{dist}(p, x), \quad
\mathcal{H}_d(M) = \{ u \in C^\infty(M) : \Delta u = 0,\ |u(x)| \le C(1+r(x))^d \text{ for some } C>0 \}.
\]
These are the harmonic functions on \(M\) of polynomial growth order at most \(d\).

It is known (Colding–Minicozzi) that for manifolds with nonnegative Ricci curvature and maximal volume growth, \(\dim \mathcal{H}_d(M)\) is finite for each \(d\) and grows at most polynomially in \(d\). For \(\mathbb{R}^n\),
\[
\dim \mathcal{H}_d(\mathbb{R}^n) = \binom{n+d}{d} - \binom{n+d-2}{d-2}.
\]

**Open Problem (Asymptotic Rigidity for Polynomial Growth Harmonic Functions):**  
Assume \( (M^n,g) \) has nonnegative sectional curvature and maximal (i.e. Euclidean) volume growth:
\[
\lim_{r\to\infty} \frac{\mathrm{Vol}(B_r(p))}{\omega_n r^n} = 1.
\]
Suppose further that for every tangent cone at infinity \( C(X) \) of \(M\), the cross-section \(X\) is a smooth Riemannian manifold with \( \mathrm{Ric}_X \ge (n-2)g_X \).

Is it true that the **asymptotic agreement** of polynomial growth harmonic function spaces,
\[
\lim_{d\to\infty} \frac{\dim \mathcal{H}_d(M)}{\dim \mathcal{H}_d(\mathbb{R}^n)} = 1,
\]
forces \( M \) to be isometric to \( \mathbb{R}^n \)?  

Equivalently: does the “asymptotic spectral completeness” of polynomial growth harmonic functions rigidly characterize Euclidean geometry among manifolds of nonnegative sectional curvature?

---

Why is it a "good" problem:  
This question sits at the intersection of **harmonic analysis**, **geometric analysis**, and **asymptotic geometry**. It proposes a new **rigidity criterion** formulated purely through the asymptotics of the spaces of harmonic functions, an analytic object that indirectly encodes the geometry at infinity.  

- **Profound connection between analysis and geometry:**  
  The dimension growth of \(\mathcal{H}_d(M)\) reflects the spectral structure of cross-sections of tangent cones at infinity. Showing that the asymptotic ratio being \(1\) forces \(M\) to be Euclidean would unify harmonic function analysis and curvature rigidity theories à la Cheeger–Colding and Colding–Minicozzi.  

- **Novel conceptual direction:**  
  Traditional rigidity results require equality in geometric quantities such as volume growth or Bishop–Gromov ratios. This problem asks whether *analytic equality*—in the asymptotic “counting measure” of harmonic functions—suffices to recover the full geometric rigidity.  

- **Technically and conceptually deep:**  
  Resolving this would require understanding how the distribution of polynomial growth harmonic functions reflects the geometry of tangent cones, possibly through a Weyl-type asymptotic linking eigenvalues of \(X\) to harmonic growth on \(M\).  

- **Plausibly open and nontrivial:**  
  Existing results give dimension estimates and rigidity when exact equality holds for all finite \(d\), but the large-\(d\) asymptotic equality case has not been established. It therefore suggests a realistic extension of known theorems beyond current techniques, opening potential connections with metric-measure spectral geometry.  

The statement is clean, conceptually motivated by known theories, and bridges distinct strands of modern differential geometry — embodying precisely the “narrow entrance, profound depth” characteristic of a good research problem.